% Load saved quantizer plots (defines \quantizerexample and \quantizerexampleste)
\newsavebox{\quantizerexample}

\begin{lrbox}{\quantizerexample}
\begin{tikzpicture}
        % Define the axis environment
    \draw[-{Stealth}] (-4.5,0) -- (4.5,0); % x-axis
    \draw[-{Stealth}] (0,-4.5) -- (0,4.5); % y-axis

    % x-axis ticks and labels
    \foreach \x in {-4, -3, -2, -1, 1, 2, 3, 4}
        \draw (\x,0.1) -- (\x,-0.1) node[below] {$\x$};
    \foreach \x in {-4, -3, -2, -1, 0,  1, 2, 3, 4}
        \fill (\x,0) circle (1.5pt); % Dots on x-axis

    % y-axis ticks and labels
    \foreach \y in {-4, -3, -2, -1, 1, 2, 3, 4}
        \draw (0.1,\y) -- (-0.1,\y) node[left] {$\y$};
    \foreach \y in {-4, -3, -2, -1, 0, 1, 2, 3, 4}
        \fill (0,\y) circle (1.5pt); % Dots on y-axis

    % Quantizer function plot (green stepped line)
    \draw[green, thick] (-4,-4) -- (-3,-4)
                        node[pos=0.5, above=2pt, black] {} % Optional: label segments
                        -- (-3,-4) -- (-3,-3)
                        -- (-3,-3) -- (-2,-3)
                        -- (-2,-3) -- (-2,-2)
                        -- (-2,-2) -- (-1,-2)
                        -- (-1,-2) -- (-1,-1)
                        -- (-1,-1) -- (0,-1)
                        -- (0,-1) -- (0,0)
                        -- (0,0) -- (1,0)
                        -- (1,0) -- (1,1)
                        -- (1,1) -- (2,1)
                        -- (2,1) -- (2,2)
                        -- (2,2) -- (3,2)
                        -- (3,2) -- (3,3)
                        -- (3,3) -- (4,3);


\end{tikzpicture}
\end{lrbox}

\newsavebox{\quantizerexampleste}

\begin{lrbox}{\quantizerexampleste}
\begin{tikzpicture}
        % Define the axis environment
    \draw[-{Stealth}] (-4.5,0) -- (4.5,0); % x-axis
    \draw[-{Stealth}] (0,-4.5) -- (0,4.5); % y-axis

    % x-axis ticks and labels
    \foreach \x in {-4, -3, -2, -1, 1, 2, 3, 4}
        \draw (\x,0.1) -- (\x,-0.1) node[below] {$\x$};
    \foreach \x in {-4, -3, -2, -1, 0,  1, 2, 3, 4}
        \fill (\x,0) circle (1.5pt); % Dots on x-axis

    % y-axis ticks and labels
    \foreach \y in {-4, -3, -2, -1, 1, 2, 3, 4}
        \draw (0.1,\y) -- (-0.1,\y) node[left] {$\y$};
    \foreach \y in {-4, -3, -2, -1, 0, 1, 2, 3, 4}
        \fill (0,\y) circle (1.5pt); % Dots on y-axis

    % Straight-Through Estimator function plot (green straight line)
    \draw[red, thick, dashed] (-4,-4) -- (4,4);


\end{tikzpicture}
\end{lrbox}


\begin{tikzpicture}[
    tensor/.style={rectangle, draw, fill=white, text width=2em, text centered, minimum height=2em},
    % % --- NUEVOS ESTILOS PARA ETIQUETAS ---
    lbl_fwd/.style={above, pos=0.5, font=\footnotesize}, % 'pos=0.5' centra la etiqueta
    lbl_bwd/.style={below, pos=0.5, font=\footnotesize},
    lbl_left/.style={left, pos=0.5, font=\footnotesize},
    lbl_right/.style={right, pos=0.5, font=\footnotesize},
    layer_function/.style={rectangle, draw, fill=white, text width=4em, text centered, rounded corners, minimum height=2em},
    operation/.style={circle, draw, fill=white, minimum size=0.1em, font=\tiny},
    arrow/.style={-latex, thick}
]

\tikzset{node distance=1.75cm and 1.5cm}

%%%% FORWARD %%%%
% draw linear block z()
\node[layer_function] (z) at (0,0) {$z()$};

% --- MODIFIED: Changed from (-2,0) to (-3,0) ---
\draw[arrow] ($(z)+(-6,0)$) -- node[lbl_fwd, pos=0.15] {$a_{m-1}$} (z.west);

% draw activation layer_function phi()
\node[layer_function, right=of z] (phi) {$\phi()$};
\node[inner sep=0pt, right=of phi] (qphi) {\scalebox{0.2}{\usebox{\quantizerexample}}};
% overlay label in the top-left corner of the quantizer plot
\node[font=\footnotesize, anchor=north west, xshift=1mm, yshift=-1mm] at (qphi.north west) {$q_{\phi}()$};
\draw[arrow] (phi) -- (qphi.west);

% draw output activation a_m (Already 3cm, no change needed)
\draw[arrow] (qphi.east) -- node[lbl_fwd, pos=0.75] {$a_{m}$} ($(qphi)+(3,0)$);

% draw arrow from z to phi
% annotate it as $z_m$
\draw[arrow] (z.east) -- node[lbl_fwd] {$z_{m}$} (phi.west);

% draw theta parameters: place them below-left of z instead of directly below
\node[inner sep=0pt, below left=of z] (qtheta) {\scalebox{0.2}{\usebox{\quantizerexample}}};
% overlay label for q_theta placed at the top-left corner
\node[font=\footnotesize, anchor=north west, xshift=1mm, yshift=-1mm] at (qtheta.north west) {$q_{\theta}()$};
\node[tensor, below=of qtheta] (theta) {$\theta$};

% theta to z
\draw[arrow] (theta) -- (qtheta);
\draw[arrow] (qtheta) -| (z);


% % --- DEFINE PARAMETER UPDATE CHAIN FIRST --- % %
\def\updateopspacing{0.5cm}
\node[operation, below=\updateopspacing of theta] (mult_eta) {$\times$};
\node[right=0.5cm of mult_eta] (eta) {$-\eta$};


%%%% BACKWARD (NOW POSITIONED RELATIVE TO UPDATE CHAIN) %%%%

% Place gradient of q_theta just below theta (use STE plot)
\node[inner sep=0pt, below=of theta] (dqtheta) {\scalebox{0.2}{\usebox{\quantizerexampleste}}};
% overlay label for dq_theta (STE) placed at the top-left corner
\node[font=\footnotesize, anchor=north west, xshift=1mm, yshift=-1mm] at (dqtheta.north west) {$\nabla_{q_{\theta}}$};

% draw z gradient to the right of dqtheta so it lines up with dqtheta
\node[layer_function, right=of dqtheta] (dz) {$\nabla_z$};

% draw phi gradiente to the right of dz
\node[layer_function, right=of dz] (dphi) {$\nabla_\phi$};

% draw q_phi gradient to the right of dphi
\node[inner sep=0pt, right=of dphi] (dqphi) {\scalebox{0.2}{\usebox{\quantizerexampleste}}};
% overlay label for dq_phi (STE) placed at the top-left corner
\node[font=\footnotesize, anchor=north west, xshift=1mm, yshift=-1mm] at (dqphi.north west) {$\nabla_{q_{\phi}}$};

% Theres a multication between upstream and dphi (below dphi)
\node[operation, below=of dphi] (mult) {$\times$};
\draw[arrow] (dphi) -- (mult);

\node[operation, below=of dqphi] (mult_qphi) {$\times$};
\draw[arrow] (dqphi) -- (mult_qphi);
\draw[arrow] ($(mult_qphi)+(3,0)$) -- (mult_qphi.east) node[lbl_bwd,pos=0.25] {$\frac{\partial \mathcal{L}}{\partial a_{m}}$};
\draw[arrow] (mult_qphi) -- (mult);


% draw a circle with $\times in the midle of dz and dphi (below dz)
\node[operation, below=of dz] (mult2) {$\times$};
\draw[arrow] (mult) -- (mult2);
\draw[arrow] (dz) -- (mult2);
% \draw[arrow] (mult2.west) -- (dz.east);

% --- MODIFIED: Changed from (-2,0) to (-3,0) ---
\draw[arrow] (mult2) -- node[lbl_bwd, pos=0.85] {$\frac{\partial \mathcal{L}}{\partial a_{m-1}}$} ($(mult2)+(-6,0)$);

% % --- CONNECT ALL THE NODES --- % %
\draw[arrow] (eta) -- (mult_eta);
\draw[arrow] (mult_eta.north) -- (theta.south);
% \draw[arrow] (negative.north) -- (mult_eta.south);

% --- MODIFIED: Added name (dtheta) to the label node ---
% start the arrow directly under dqtheta by using the vertical projection
\draw[arrow] (dqtheta |- mult2) -- node[lbl_left] (dtheta) {$\frac{\partial \mathcal{L}}{\partial \theta}$} (dqtheta);

\draw[arrow] (dqtheta) -- (mult_eta);

% % --- BACKGROUND LAYER BOX --- % %
\begin{pgfonlayer}{background}
    \node [
        draw,
        dotted,
        thick,
        orange,
        rounded corners,
        % --- MODIFIED: Added (dtheta) to fit list ---
        fit=(z) (phi) (theta) (mult_eta) (dz) (dphi) (mult) (mult2) (dtheta) (qphi) (dqphi) (dqtheta) (mult_qphi),
        % --- MODIFIED: Reduced inner sep for better spacing ---
        inner sep=0.25cm
    ] (layer_box) {};
    % Add horizontal line through the box at theta level (background so it sits behind nodes)
    \draw [gray, dashed, thin] ($(theta)+(-3,0)$) -- ($(theta)+(13,0)$);
\end{pgfonlayer}

% --- Placed your label relative to the box's corner ---
\node[font=\footnotesize, orange, above right, xshift=2mm] at (layer_box.north east) {$L_m$};


\end{tikzpicture}
